% =============================================================================
% MACROS.TEX — Macros propios (complementan al paquete physics)
% =============================================================================
%
% El paquete physics ya provee:
%   \abs, \norm, \tr, \grad, \div (divergencia), \curl, \laplacian,
%   \dv, \pdv, \bra, \ket, \braket, \ip (producto interno),
%   \op (producto exterior), \comm (conmutador), \acomm,
%   \rank, \erf, \Res, \Re, \Im, \pv, etc.
%
% Aquí solo definimos lo que physics NO cubre.
% =============================================================================

% ─── Conjuntos numéricos ─────────────────────────────────────────────────────
\newcommand{\R}{\mathbb{R}}
\newcommand{\N}{\mathbb{N}}
\newcommand{\Z}{\mathbb{Z}}
\newcommand{\C}{\mathbb{C}}
\newcommand{\Q}{\mathbb{Q}}

% ─── Transpuesta ─────────────────────────────────────────────────────────────
\newcommand{\T}{^\top}

% ─── Optimización ────────────────────────────────────────────────────────────
\DeclareMathOperator*{\argmin}{arg\,min}
\DeclareMathOperator*{\argmax}{arg\,max}

% ─── Aprendizaje automático / estadística ────────────────────────────────────
\newcommand{\Loss}{\mathcal{L}}
\newcommand{\D}{\mathcal{D}}
\newcommand{\Params}{\Theta}
\newcommand{\params}{\theta}
\DeclareMathOperator{\KL}{KL}
\DeclareMathOperator{\Var}{Var}
\DeclareMathOperator{\Cov}{Cov}
\newcommand{\E}[1]{\mathbb{E}\!\left[#1\right]}

% ─── Física: notación específica ─────────────────────────────────────────────
\newcommand{\Lap}{\laplacian}                   % alias → \laplacian de physics
\newcommand{\Div}{\divergence}                  % alias → \divergence de physics
\newcommand{\keff}{k_{\mathrm{eff}}}
\newcommand{\score}{\grad_{\bm{x}}\log p}
\newcommand{\Ep}[1]{\langle #1 \rangle}
\newcommand{\FP}{\mathcal{L}_{\mathrm{FP}}}

% ─── Ecuaciones diferenciales estocásticas ───────────────────────────────────
\newcommand{\SDE}[3]{d#1 = #2\,dt + #3\,dW}

% ─── Símbolos auxiliares ─────────────────────────────────────────────────────
\newcommand{\OK}{\ensuremath{\checkmark}}

% =============================================================================
% NOTA: si necesitas \rank como \DeclareMathOperator{rango} en lugar de
% la versión de physics, descomenta:
% \let\rank\relax
% \DeclareMathOperator{\rank}{rango}
% =============================================================================